% Fig. 2: the box comparison of Appendix D (redrawn for the revised edition of 7 November 2026;
% it replaces "two affine paths, one quantifier", which drew the reading of the retired item
% A8(vi)).  Left: the GJS squares Delta_{l,k} (row l in the time x^0, column k in the space
% x^1, centred at integer points), the observable's squares, the box B and one separating
% strip of each kind: (a) Theorem D.3 (box B_s(g,I): rows l_s-1..l_s+1 times every column of
% supp g and I; every charged square outside lies above or below; parallel strip P_l of width
% |l-l_s|-1), drawing values l_s = 0, columns of supp g -3..2, I in columns -1..0, slab rows
% -3..3; (b) Theorem D.4 (S_a = row 0, columns 0..2; B_{3x5} = rows -1..1, columns -1..3; ring j
% = squares with delta = j; class (P) |l| = 1+j, parallel strip of width j; class (T) |l| <= j
% in the side columns, transverse strip of width j = A8(iv)-T), drawing values: rings 1-2, a
% cutoff h drawn as an irregular region, the path h_alpha = h on B and alpha h outside B
% (Lemma N.5(ii)).  Right: Lemma D.1, Lemma D.2, the box term (Lemma C.2, Theorem C.5) and the
% two theorems.  Drawn from the statements; nothing is computed; no page image.
\begin{tikzpicture}[font=\scriptsize,
  sq/.style={draw=gline,line width=0.3pt},
  lbox/.style={draw=mut,rounded corners=1pt,fill=white,align=left,inner sep=3pt,text width=6.0cm},
  tlbox/.style={lbox,draw=warn,fill=tfill},
  arr/.style={-{Latex[length=1.5mm]},line width=0.5pt,draw=mut}]
  % ================= (a) Theorem D.3
  \begin{scope}[x=0.37cm,y=0.37cm]
    % slab h_{T,g}: columns -3..2, rows -3..3
    \fill[acc!9!white] (-3.5,-3.5) rectangle (2.5,3.5);
    % parallel strip P_3 for the squares of row 3: 0.5 < x^0 < 2.5
    \fill[pattern=north east lines,pattern color=acc!55!white] (-4.5,0.5) rectangle (3.5,2.5);
    \foreach \k in {-4,...,3} \foreach \l in {-3,...,3}
      \draw[sq] (\k-0.5,\l-0.5) rectangle (\k+0.5,\l+0.5);
    % squares of the observable, row 0, columns -1..0
    \fill[acc] (-1.5,-0.5) rectangle (0.5,0.5);
    % box B_s(g,I): rows -1..1, columns -3..2
    \draw[ok,line width=1.1pt] (-3.5,-1.5) rectangle (2.5,1.5);
    % one charged square outside, row 3
    \draw[acc,line width=0.9pt] (0.5,2.5) rectangle (1.5,3.5);
    \draw[gline,->] (-4.9,0) -- (3.9,0) node[right,mut,inner sep=1pt] {$x^1$};
    \draw[gline,->] (-4.5,-3.9) -- (-4.5,3.9) node[above,mut,inner sep=1pt] {$x^0$};
    \node[font=\footnotesize\bfseries,anchor=south west,inner sep=0pt] at (-4.9,4.25) {(a)};
    \node[align=left,anchor=north west,text width=3.3cm,inner sep=0pt] at (-4.6,-3.85)
      {Thm D.3: the box $B$ (green) spans every column of $\supp g$; every square outside lies above or below, so the strip (hatched) is parallel};
  \end{scope}
  % ================= (b) Theorem D.4
  \begin{scope}[shift={(3.85cm,0)},x=0.37cm,y=0.37cm]
    % a cutoff h in C_E: an irregular region; alpha h outside B is drawn lighter
    \fill[acc!9!white] (-3.4,-2.9) -- (-1.8,-3.4) -- (2.6,-3.3) -- (5.4,-1.6) -- (5.3,2.3) -- (3.2,3.4) -- (-1.2,3.2) -- (-3.4,1.6) -- cycle;
    \fill[acc!22!white] (-1.5,-1.5) rectangle (3.5,1.5);
    % transverse strip for the class-(T) square in row 0, column 4: 2.5 < x^1 < 3.5
    \fill[pattern=north west lines,pattern color=warn!70!white] (2.5,-3.5) rectangle (3.5,3.5);
    % parallel strip for the class-(P) square in row 2: 0.5 < x^0 < 1.5
    \fill[pattern=north east lines,pattern color=acc!55!white] (-3.5,0.5) rectangle (5.5,1.5);
    \foreach \k in {-3,...,5} \foreach \l in {-3,...,3}
      \draw[sq] (\k-0.5,\l-0.5) rectangle (\k+0.5,\l+0.5);
    % ring 1 (dashed outline)
    \draw[mut,dashed,line width=0.5pt] (-2.5,-2.5) rectangle (4.5,2.5);
    % squares of the observable S_a: row 0, columns 0..2
    \fill[acc] (-0.5,-0.5) rectangle (2.5,0.5);
    % the box B_{3x5}
    \draw[ok,line width=1.1pt] (-1.5,-1.5) rectangle (3.5,1.5);
    % one class-(P) square (row 2) and one class-(T) square (row 0, column 4)
    \draw[acc,line width=0.9pt] (0.5,1.5) rectangle (1.5,2.5);
    \fill[tfill] (3.5,-0.5) rectangle (4.5,0.5);
    \draw[warn,line width=0.9pt] (3.5,-0.5) rectangle (4.5,0.5);
    \node[font=\tiny,inner sep=0.5pt] at (1,2) {P};
    \node[font=\tiny,inner sep=0.5pt,warn] at (4,0) {T};
    \node[font=\tiny,inner sep=0.5pt,mut,fill=white] at (-2.6,-3.0) {$\alpha h$};
    \node[font=\tiny,inner sep=0.5pt,fill=white] at (1,-1.0) {$h$};
    \node[font=\footnotesize\bfseries,anchor=south west,inner sep=0pt] at (-3.9,4.25) {(b)};
    \node[align=left,anchor=north west,text width=3.5cm,inner sep=0pt] at (-3.6,-3.85)
      {Thm D.4: the box $B_{3\times5}$ and its innermost ring (dashed); a square in the top or bottom row of a ring (P) gets a parallel strip, a side square (T) a transverse one: A8(iv)-T};
  \end{scope}
  % ================= right: the steps
  \node[lbox,anchor=north west] (s1) at (6.45cm,1.65cm)
    {(1) Lemma D.1: along $h_\alpha=h^B+\alpha(h-h^B)\in\CE$, the difference $\langle F\rangle_h-\langle F\rangle_{h^B}$ is\\ $-\int_0^1\mathrm{Cov}_{h_\alpha}(F,\mathcal U_{h-h^B})\,d\alpha$};
  \node[lbox,anchor=north west] (s2) at ($(s1.south west)+(0,-0.12)$)
    {(2) Lemma D.2: $\mathcal U_{h-h^B}$ is a sum over the squares $\Delta$ outside $B$; Thm 1.1.7 across a strip of width $w_\Delta$ bounds each by $C_7c_U\|F\|_{\rm GJS}e^{-m_1w_\Delta}$};
  \node[lbox,anchor=north west] (s3) at ($(s2.south west)+(0,-0.12)$)
    {(3) the box term $\langle F\rangle_{h^B}$: Lemma C.2 bounds $Z(h^B)$ below, Thm C.5 makes $S_2(B)$ finite; no clustering};
  \node[lbox,anchor=north west,text width=2.8cm] (o3) at ($(s3.south west)+(0,-0.45)$)
    {Thm D.3: fixed $g$, uniform in $T$; parallel strips only $\to$ \S7};
  \node[tlbox,anchor=north east,text width=2.8cm] (o4) at ($(s3.south east)+(0,-0.45)$)
    {Thm D.4: every $h\in\CE$; transverse strips too $\to$ Thm 5.1};
  \draw[arr] (s3.south -| o3) -- (o3.north);
  \draw[arr] (s3.south -| o4) -- (o4.north);
\end{tikzpicture}
